\documentclass[tikz,border=12pt]{standalone}
\usepackage[utf8]{inputenc}
\usepackage[T1]{fontenc}
\usepackage{lmodern}
\usepackage{amsmath,amssymb}
\usepackage{xcolor}

\usetikzlibrary{shapes.geometric,arrows.meta,positioning,fit,calc,shadows.blur}

% Academic Color Palette
\definecolor{k1blue}{RGB}{37, 99, 235}
\definecolor{k1bg}{RGB}{239, 246, 255}
\definecolor{k1border}{RGB}{147, 197, 253}

\definecolor{k2purple}{RGB}{124, 58, 237}
\definecolor{k2bg}{RGB}{245, 243, 255}
\definecolor{k2border}{RGB}{196, 181, 253}

\definecolor{consgreen}{RGB}{22, 163, 74}
\definecolor{consbg}{RGB}{240, 253, 244}
\definecolor{consborder}{RGB}{134, 239, 172}

\definecolor{k3orange}{RGB}{234, 88, 12}
\definecolor{k3bg}{RGB}{255, 247, 237}
\definecolor{k3border}{RGB}{253, 186, 116}

\definecolor{darkslate}{RGB}{15, 23, 42}
\definecolor{lightgray}{RGB}{241, 245, 249}

\begin{document}
\sffamily

\begin{tikzpicture}[
    >=Stealth,
    node distance=1.2cm and 1.8cm,
    font=\small,
    box/.style={
        rectangle, rounded corners=6pt, draw=#1, line width=1.2pt,
        align=center, inner sep=8pt, blur shadow={shadow blur steps=4, shadow opacity=15}
    },
    subbox/.style={
        rectangle, rounded corners=4pt, draw=#1, line width=0.8pt, fill=white,
        align=center, inner sep=6pt
    },
    diamondbox/.style={
        diamond, draw=darkslate, line width=1.5pt, fill=white, aspect=1.8,
        align=center, inner sep=4pt, blur shadow={shadow blur steps=4, shadow opacity=15}
    },
    line/.style={
        draw=#1, line width=1.5pt, ->
    }
]

% 1. INPUT NODE
\node[box=darkslate, fill=white, text width=9.5cm] (input) {
    \textbf{\large GİRDİ (DOĞRULAMA TALEBİ)}\\[3pt]
    \footnotesize
    \textbf{Bağlam Metni ($C$):} Resmi Kanıt Dokümanı / Mevzuat / Klinik Kılavuz\\
    \textbf{Doğrulanacak İddia ($S$):} LLM Yanıtı (Tıp, Hukuk, Finans)
};

% 2. COMPONENT 1 (K1 - HOLISTIC ENCODER)
\node[box=k1blue, fill=k1bg, text width=5.6cm, below left=1.4cm and 0.2cm of input] (k1) {
    \textbf{\textcolor{k1blue}{\normalsize BİLEŞEN 1: BÜTÜNCÜL ENCODER (K1)}}\\[5pt]
    \begin{minipage}{5.4cm}
        \centering
        \begin{tikzpicture}[font=\footnotesize]
            \node[subbox=k1border, text width=5.0cm] (k1_model) {
                \textbf{dbmdz/electra-base-turkish-cased}\\
                (110M Parametre, All-Linear LoRA)
            };
            \node[subbox=k1border, text width=5.0cm, below=0.35cm of k1_model] (k1_inf) {
                \textbf{Küresel Sekans Sınıflandırması}\\
                $\text{Girdi: } [C] \oplus [S] \to \text{Softmax}$\\
                Gecikme: $\approx 12$\,ms
            };
            \draw[->, line width=1pt, draw=k1blue] (k1_model) -- (k1_inf);
        \end{tikzpicture}
    \end{minipage}\\[6pt]
    \textbf{\textcolor{k1blue}{K1 Kararı: $\hat{y}_{\text{K1}} \in \{0, 1, 2, 3\}$}}
};

% 3. COMPONENT 2 (K2 - PROPOSITION NLI)
\node[box=k2purple, fill=k2bg, text width=6.2cm, below right=1.4cm and 0.2cm of input] (k2) {
    \textbf{\textcolor{k2purple}{\normalsize BİLEŞEN 2: ÖNERME-DÜZEYLİ NLI (K2)}}\\[5pt]
    \begin{minipage}{6.0cm}
        \centering
        \begin{tikzpicture}[font=\footnotesize]
            \node[subbox=k2border, text width=5.6cm] (k2_atom) {
                \textbf{1. Aşama: Gemma-4 QLoRA Atomizer}\\
                $\text{İddia } S \to \{a_1, a_2, \dots, a_n\}$ Önermeler\\
                \scriptsize \%99.58 Test Kapsaması (Ort. 2.08 atom/iddia)
            };
            \node[subbox=k2border, text width=5.6cm, below=0.35cm of k2_atom] (k2_nli) {
                \textbf{2. Aşama: mDeBERTa-v3 NLI (278M)}\\
                $(C, a_i) \to \{p_{\text{ent}}, p_{\text{neu}}, p_{\text{con}}\}$\\
                Soft-Probability Toplulaştırma ($\approx 28$\,ms)
            };
            \draw[->, line width=1pt, draw=k2purple] (k2_atom) -- (k2_nli);
        \end{tikzpicture}
    \end{minipage}\\[6pt]
    \textbf{\textcolor{k2purple}{K2 Kararı: $\hat{y}_{\text{K2}} \in \{0, 1, 2, 3\}$}}
};

% ARROWS FROM INPUT TO K1 & K2
\draw[line=k1blue] (input.south) -| node[above, pos=0.25, font=\scriptsize, text=k1blue] {} (k1.north);
\draw[line=k2purple] (input.south) -| node[above, pos=0.25, font=\scriptsize, text=k2purple] {} (k2.north);

% 4. DECISION GATE (ROUTER)
\coordinate (midk) at ($(k1.south)!0.5!(k2.south)$);
\node[diamondbox, below=1.0cm of midk] (gate) {
    \textbf{Kararlar}\\
    \textbf{Aynı mı?}\\
    \scriptsize $(\hat{y}_{\text{K1}} \stackrel{?}{=} \hat{y}_{\text{K2}})$
};

% ARROWS FROM K1 & K2 TO GATE
\draw[line=k1blue] (k1.south) |- (gate.west);
\draw[line=k2purple] (k2.south) |- (gate.east);

% 5. BRANCH A: CONSENSUS NODE
\node[box=consgreen, fill=consbg, text width=5.6cm, below left=1.2cm and 0.4cm of gate] (consensus) {
    \textbf{\textcolor{consgreen}{\large \checkmark\ TAM KONSENSÜS BÖLGESİ}}\\[4pt]
    \footnotesize
    \textbf{Kapsam:} 358 Altın Test Örneği (\%74.90)\\[2pt]
    \textbf{Doğruluk:} \textbf{\textcolor{consgreen}{\%94.13}} (337/358 Doğru)\\[2pt]
    \textbf{Maliyet / Hız:} 0 LLM Çağrısı (\$0.00), $\approx 35$\,ms\\[3pt]
    \scriptsize İki bağımsız mimari doğrudan uzlaştığı için\\karar dış müdahale olmadan onaylanır.
};

% 6. BRANCH B: COMPONENT 3 (K3 - META-JUDGE)
\node[box=k3orange, fill=k3bg, text width=6.2cm, below right=1.2cm and 0.4cm of gate] (k3) {
    \textbf{\textcolor{k3orange}{\normalsize BİLEŞEN 3: EMSALLİ META-HAKEM (K3)}}\\[4pt]
    \footnotesize
    \textbf{Hakem Modeli:} Meta Llama-3.3-70B-Instruct\\
    \scriptsize (Ablasyon: Qwen-2.5-72B \%70.0 $\mid$ GPT-4o-mini \%63.3)\\[3pt]
    \begin{minipage}{6.0cm}
        \centering
        \begin{tikzpicture}[font=\scriptsize]
            \node[subbox=k3border, text width=5.6cm] (k3_input) {
                \textbf{Hakeme Sunulan Dosya:}\\
                Bağlam $C$ + İddia $S$ + K1 Görüşü + K2 Atom Raporu\\
                + 4 Gerçek Emsal Mahkeme İçtihadı (Few-Shot)
            };
        \end{tikzpicture}
    \end{minipage}\\[4pt]
    \textbf{Uyuşmazlık Çözüm Başarımı:} \textbf{\textcolor{k3orange}{\%75.83}} (91/120 vaka)\\
    \scriptsize 120 Vakanın Toplam Maliyeti: Yalnızca \$0.036 USD
};

% GATE BRANCHING ARROWS
\draw[line=consgreen] (gate.south) -| node[above right, pos=0.15, font=\bfseries\footnotesize, text=consgreen] {EVET (\%74.9)} (consensus.north);
\draw[line=k3orange] (gate.south) -| node[above left, pos=0.15, font=\bfseries\footnotesize, text=k3orange] {HAYIR (\%25.1)} (k3.north);

% 7. FINAL OUTPUT NODE
\coordinate (midbottom) at ($(consensus.south)!0.5!(k3.south)$);
\node[box=darkslate, fill=darkslate, text=white, text width=11.2cm, below=1.2cm of midbottom] (output) {
    \textbf{\large \textcolor{cyan}{NİHAİ HİBRİT DOĞRULAMA ÇIKTISI $\hat{y}_{\text{Hibrit}}$}}\\[4pt]
    \normalsize
    \textbf{478 Altın Test Kümesi Başarımı: \textcolor{green!40}{\%89.54 Doğruluk} $\;\vert\;$ \textcolor{green!40}{Macro-F1: 0.8966} $\;\vert\;$ \textcolor{green!40}{MCC: 0.8628}}\\[3pt]
    \footnotesize
    \textcolor{gray!20}{McNemar İstatistiki Anlamlılık: vs K1 ($p = 0.00064$) $\;\vert\;$ vs K2 ($p = 1.34 \times 10^{-12}$)}
};

% CONNECT TO FINAL
\draw[line=consgreen] (consensus.south) |- (output.west);
\draw[line=k3orange] (k3.south) |- (output.east);

\end{tikzpicture}

\end{document}
